% ── §8: Binary Sequence Spaces and Their Topological Structure ──
\section{Binary Sequence Spaces and Their Topological Structure}
\label{sec:binary-spaces}

A central theoretical contribution of this work is the recognition that the
Gray-code encoding of biological alphabets induces a metric structure on
sequence space that is simultaneously compatible with the Karnaugh-map grid
and reflective of biophysical similarity.  We develop this idea formally.

\subsection{The Encoded Sequence Space}

Let $\Sigma = \{a_0, \ldots, a_{19}\}$ denote the canonical amino-acid
alphabet, and let $\phi : \Sigma \to \{0, \ldots, 19\}$ assign each residue
its He~2012 integer code.  The reflected binary Gray transformation $g(n) =
	n \oplus (n \gg 1)$ maps these codes into $\{0, \ldots, 31\} \subset \{0,1\}^5$,
embedding the alphabet into the five-dimensional hypercube $Q_5$ whose
vertices are the 32 five-bit vectors.  Twenty of these vertices are occupied;
the remaining twelve constitute the don't-care region that gives the
Quine--McCluskey algorithm its freedom to form larger implicants.

For two residues $a, b \in \Sigma$, define their \emph{encoding distance} as
\begin{equation}
	d_\Sigma(a, b) = d_H\!\bigl(g(\phi(a)),\, g(\phi(b))\bigr),
	\label{eq:encoding-distance}
\end{equation}
where $d_H$ is the Hamming metric on $\{0,1\}^5$.  This function satisfies
the metric axioms (non-negativity, symmetry, triangle inequality) because it
is the pullback of a metric under an injective map.  It defines a graph
structure on $\Sigma$ in which two residues are adjacent iff $d_\Sigma(a,b)
	= 1$; this adjacency graph is a subgraph of $Q_5$ with exactly 40 edges
(ordered count: 80), as verified by exhaustive enumeration in
Lean~4~(Theorem~\texttt{cc\_he\_dist1\_ordered\_eq\_80}).

\subsection{Contact Enrichment in the Adjacency Graph}

The central structural hypothesis of our framework is that pairs of residues
connected by an edge in the encoding-adjacency graph are enriched among
native contacts relative to what would be expected from their marginal
frequencies alone.  This hypothesis was confirmed on the PSICOV150 benchmark
with overwhelming statistical significance ($p < 10^{-77}$,
Section~\ref{sec:results-gray}) and further validated by a permutation test
showing that the specific He/Gray assignment outperforms random code
assignments ($z = 2.51$).

We can understand this enrichment through the lens of \emph{information
	geometry}.  The mutual information $\MI(i,j)$ between positions $i$ and $j$
measures how much knowing the residue at one position reduces uncertainty
about the other.  When both positions are buried in the protein core, their
residues are drawn from a restricted alphabet dominated by hydrophobic
members of the same physicochemical group.  Under the He/Gray encoding,
members of the same group occupy nearby codewords (within-group mean distance
$\leq 2$ bits), so contacting pairs of core residues tend to have small
$d_\Sigma$, producing the observed enrichment.

\subsection{Karnaugh Maps as Quotient Structures}

A Karnaugh map of dimension $(w_1, w_2)$ over encoded alphabets $\Sigma_1$
and $\Sigma_2$ is the Cartesian product $Q_{w_1} \times Q_{w_2}$ arranged as
a $2^{w_1} \times 2^{w_2}$ grid with Gray-code row and column ordering.  For
protein-pair analysis, $w_1 = w_2 = 5$, giving a $32 \times 32$ grid with
1024 cells.  The ternary labeling (On/Off/Don't-care) partitions these cells
into three disjoint sets whose semantics are fixed by the data:
\begin{align}
	\mathrm{ON}(i,j)  & = \{(a,b) : f_{ij}(a,b) > 0 \wedge (a,b) \neq
	(\pi_i, \pi_j)\},                                                 \\
	\mathrm{OFF}(i,j) & = \{(a,b) : f_{ij}(a,b) = 0 \wedge (a,b) \neq
	(\pi_i, \pi_j)\},                                                 \\
	\mathrm{DC}(i,j)  & = \{(a,b) : (a,b) = (\pi_i, \pi_j)\}
	\cup \{a \geq 20\} \cup \{b \geq 20\},
\end{align}
where $f_{ij}(a,b)$ is the joint frequency and $(\pi_i, \pi_j)$ is the
reference pair.

Quine--McCluskey minimisation operates on this partition by finding the
coarsest sum-of-products expression whose implicants cover all ON cells while
avoiding all OFF cells.  Don't-care cells serve as free variables that may be
included in or excluded from any implicant without affecting correctness.
The result is a set of prime implicants, each representing a maximal
hypercube of the Karnaugh map that lies entirely within the ON $\cup$ DC
region---a geometric object that corresponds to a conjunction of amino-acid
identity constraints at the two positions.

\subsection{Homotopical Interpretation}

The hypercube $Q_w$ has a well-known topological interpretation as the
$categorical product$ of $w$ copies of the unit interval $[0,1]$, making it a
contractible space with trivial homotopy groups in all dimensions $\geq 1$.
When we restrict to the twenty occupied vertices corresponding to canonical
amino acids, we obtain a subcomplex whose topology encodes the
physicochemical organisation of the genetic code.  Specifically, the
within-group adjacency structure creates connected components (one per
physicochemical class), and the between-group adjacencies create bridges
between these components that reflect shared biochemical properties.

From this perspective, a coevolutionary signal appears as a concentration of
interaction energy along particular edges and faces of this subcomplex.
Positions that coevolve tend to involve residues whose codewords lie on the
same face of the hypercube (small $d_\Sigma$), because such pairs share
physicochemical properties that allow simultaneous substitution without
disrupting the fold.  The Quine--McCluskey minimiser detects precisely these
concentrated regions by finding maximal subcubes within which all observed
combinations are ON.

\subsection{Connection to Constraint-Based Protein Design}

The SOP expressions produced by QM minimisation have a direct interpretation
as \emph{design constraints}.  Each prime implicant specifies a conjunction
of conditions (amino-acid identities at two positions, plus a separation
range) that is sufficient for contact formation.  A complete set of such
implicants defines a Boolean circuit whose ON output characterises the
allowed sequence space for a given fold.

In inverse protein folding, one seeks sequences compatible with a target
structure.  Our framework provides a natural bridge: the circuit's ON region
is exactly the subset of sequence space consistent with the observed
coevolutionary constraints, and sequences outside this region are predicted
to be incompatible.  While current precision does not suffice for de novo
folding (Section~\ref{sec:results-reconstruction}), the constraint
satisfaction perspective suggests a route to improvement through integration
with energy-based methods that operate on the same discrete representation.
